\cotefile{74}
\scanpage{56}
Sa structure combinatoire est donnée par les ens $I_{\gamma}$ de cardinal 3 et $\deleted{\unsure{\Sigma|I}}\Lambda_{\gamma}=\Lambda-\{\gamma\}\inserted{=\{\alpha,\beta\}}$ de cardinal 2, comme $H\simeq I_{\gamma}\times\Lambda_{\gamma}$, avec $\tau_{\deleted{\unsure{\alpha}}i,\alpha}$ et $\tau_{j\beta}$ \uline{liés} ssi $i\ne j,\ \alpha\ne\beta$. \deleted{\unsure{(deux îlots}}

\diagram{\begin{tikzcd}[ampersand replacement=\&,column sep=small,row sep=small]
\&\tau_{i\alpha}\arrow[r,dash]\&\tau_{j\beta}\arrow[dr,dash]\&\\
\tau_{k\beta}\arrow[ur,dash]\arrow[dr,dash]\&\&\&\tau_{k\alpha}\arrow[dl,dash]\\
\&\tau_{j\alpha}\arrow[r,dash]\&\tau_{i\beta}\&
\end{tikzcd}}

\textcircled{14} Les triangles de triades biliées correspondent aux couples $(\alpha,\beta)$ d'élt. distincts de $\Lambda$ (i.e. aux repères de $\Lambda$) en associant à tt tel repère le \deleted{\unsure{ens de}} triangle de triades $T_{\alpha,\beta,\gamma}\deleted{\unsure{\{\tau_{i,\alpha},\overline{\tau}_{j}}}\ \{\tau_{i\alpha}\}_{i\in I_{\gamma}}$. Cet ens est donc un torseur à droite sous le groupe $\unsure{\mathfrak S_3|\mathscr S_3}$. Ceci dit, \deleted{\unsure{le triangle de}} \inserted{le triangle des} triades \uline{opposé} à $T_{\alpha\beta\gamma}$ est
\[
\overline{T_{\alpha\beta\gamma}}=T_{\beta,\alpha,\gamma}
\]

Le triangle de triades $\widetilde T_{\alpha\beta\gamma}$ \deleted{\unsure{co}} contenu dans le $\unsure{=|\simeq}$ bitriangle $X_\alpha$ que $\deleted{\unsure{X_{\alpha\beta}}}T_{\alpha\beta\gamma}$, et distinct de $T_{\alpha\beta}$, est $T_{\alpha\gamma\beta}$. Donc le groupe engendré par les opérations $-$ et $\sim$ sur l'ens des triangles de triades liées

% archive sheet number: 55

